\section{Related Work}
\label{sec:related_work}

\para{Generative and edit-based linguistic steganography.} Earlier text steganography work focused on dense encoding schemes that steer generation or edit local tokens while preserving fluency. Adaptive dynamic grouping gives a formal generative baseline for secure text steganography \cite{zhang2021provably}. Edit-based methods using masked language models show that sparse local substitutions can be much harder to detect than naïve baselines \cite{ueoka2021masked}. More recent LLM-era systems improve quality and capacity with in-context generation or black-box prompting, including zero-shot generative steganography and keyword-set methods for large models \cite{lin2024zeroshot,wu2024generative}. Raj-Sankar and Rajagopalan argue that density, detectability, and decodability should be treated as linked design variables \cite{rajsankar2023science}. Our work builds on that framing, but instead of proposing a new hiding method, we isolate dilution as an evaluation axis for present-day model-based steganalysis.

\para{Steganalysis with language models.} A second thread studies whether learned language representations help detect hidden text. Yi et al. show that pretrained language-model features improve steganalysis over older RNN and CNN baselines in several settings \cite{yi2021efficient}. Bai et al. extend this idea by using LLMs directly as steganalyzers and report large gains over classical detectors on dense or structurally simpler tasks \cite{bai2024nextgen}. These results motivate our use of prompted LLMs as black-box analyzers, but they do not answer whether the same models can recover a sparse secret when evidence is scattered across long context. We focus on that harder extraction setting while still measuring binary detection.

\para{Robustness and frontier-model covert capability.} Recent work has shifted from classical steganography metrics toward robustness and safety questions. Perry et al. distinguish local and semantic robustness and show that some LLM-based schemes survive paraphrase or light edits better than watermark-like methods \cite{perry2025robust}. Karpov et al. study covert communication and hidden reasoning in frontier models, while Zolkowski et al. find that current models show only early and limited steganographic capability without stronger coordination affordances \cite{karpov2025potentials,zolkowski2025early}. Our benchmark is closest in spirit to this safety-oriented line, but it tests a different capability: scheme-aware detection and extraction from diluted keyword signals over realistic cover corpora rather than coordinated message passing or hidden reasoning tasks.

\para{Positioning our work.} Across these threads, the missing benchmark is one that explicitly sweeps signal density and separates detection from recovery. That is the niche this paper fills. We ask whether current LLMs can still recover a message once the hiding rule remains simple in principle but becomes sparse, dispersed, and confounded by distractor triggers.
